\subsection{Matriz de Evaluación Ponderada}

Con base en los criterios anteriores, se aplicó una matriz de evaluación ponderada en una
escala discreta de 1 a 5, donde 5 representa el mayor grado de adecuación al objetivo del
trabajo y 1 el menor. Los puntajes corresponden a una valoración relativa respecto a la
conveniencia de cada biblioteca para integrarse en una arquitectura de almacenamiento
comprimido por bloques dentro del simulador seleccionado.

La asignación de valores se apoyó en la documentación oficial y en la literatura técnica
asociada a cada biblioteca. Estas fuentes permitieron identificar el propósito de diseño de
cada biblioteca, su modelo de uso sobre memoria, la presencia o ausencia de mecanismos
internos de particionamiento, y su orientación hacia datos binarios o numéricos.

\begin{table}[H]
    \centering
    \caption{Evaluación de bibliotecas de compresión sin pérdida (escala 1--5) y total ponderado}
    \label{tab:evaluacion_bibliotecas}
    \begin{tabular}{lccccc}
        \hline
        \textbf{Criterio} & \textbf{Peso} & \textbf{Blosc2} & \textbf{Zstd} & \textbf{LZ4} & \textbf{zlib} \\
        \hline
        C1 & 0.35 & 5 & 3 & 3 & 2 \\
        C2 & 0.30 & 4 & 4 & 5 & 2 \\
        C3 & 0.20 & 5 & 4 & 4 & 4 \\
        C4 & 0.15 & 5 & 2 & 2 & 2 \\
        \hline
        \textbf{Total} & \textbf{1.00} & \textbf{4.75} & 3.50 & 3.65 & 2.50 \\
        \hline
    \end{tabular}
\end{table}

La matriz muestra una diferencia clara entre una biblioteca concebida con afinidad hacia
datos numéricos particionados y bibliotecas generalistas de compresión sin pérdida.
Blosc2 obtuvo la mayor valoración en adecuación a almacenamiento por bloques porque
su modelo incorpora contenedores fragmentados, filtros reversibles y operación en
memoria dentro de un marco unificado \parencite{blosc2_docs}. Zstd y LZ4 alcanzaron una
valoración favorable en costo de acceso e integración debido a sus APIs en memoria y a
su perfil de alto desempeño, pero dependen en mayor medida de que el particionamiento,
la política de acceso y la gestión de bloques sean resueltos por la aplicación
\parencite{ZstdRepo, ZstdSite, LZ4Repo}. zlib conserva valor como referencia estable y madura,
aunque su perfil resulta menos atractivo para una ruta crítica sensible a latencia y acceso
repetido sobre fragmentos del vector de estado \parencite{ZlibTechDetails}.

Las diferencias más relevantes se concentran en los criterios con mayor peso. Blosc2
sobresale en estructura por bloques e integración porque ya articula un modelo de trabajo
alineado con datos binarios y numéricos en memoria. Zstd y LZ4 ofrecen códecs rápidos
y ampliamente adoptados, pero trasladan a la aplicación la responsabilidad principal
sobre la organización del almacenamiento comprimido. zlib, aunque robusta y muy
conocida, resulta menos favorable cuando la prioridad es reducir la sobrecarga temporal
del acceso repetido a fragmentos del estado.
